\chapter{ Tokens }

  Experimentando con algunos modelos de lenguaje, se me ocurrio replicar el uso
  de tokens: usar regex pero para crear una \textbf{tabla de tokens}. No importa
  el contenido, solo el tipo, si es un comando, texto, entorno, etc. Ademas note
  que el arbol resultante no necesita pensar en como renderizar sino en como
  organizar la salida, el navegador se encarga del resto.

  Los primeros comandos integrados fueron una prueba: \backslash section,
  \backslash\backslash y estilos de texto, para verificar que la tokenizacion
  funcionaba.

  \section{ Estructura de archivos }

  Cuando el tokenizador crecio, separar en archivos de una sola responsabilidad
  fue necesario. Seguir el flujo en un unico archivo se habia vuelto dificil,
  y el parseo asincrono tenia dos puntos de ejecucion por step. Se elimino.

  \section{ Entornos }

  Con esa base mas limpia, los primeros entornos fueron \texttt{enumerate}
  e \texttt{itemize}. En retrospectiva no fue la mejor eleccion porque estos
  entornos usan \backslash item como separador, un \textit{cierre implicito}
  sin marca final. Eso complica la deteccion, pero si funciona bien la
  misma logica sirve para otros elementos similares.

  \begin{itemize}
      \item Test de lista
      \item Prueba con inline math $\int 2x^{3y}$
      \item Cierre implicito del anterior
  \end{itemize}
